% Generated draft: requires mathematical and source review.
\documentclass[11pt]{article}
\usepackage[margin=25mm]{geometry}
\usepackage{fontspec}
\setmainfont{Latin Modern Roman}
\usepackage{amsmath,amssymb,mathtools,unicode-math}
\setmathfont{Latin Modern Math}
\usepackage{enumitem,needspace,xurl,hyperref}
\setlength{\parindent}{0pt}
\setlength{\parskip}{5pt}
\setlength{\emergencystretch}{3em}
\newcommand{\N}{\mathbb{N}}
\newcommand{\Z}{\mathbb{Z}}
\newcommand{\Q}{\mathbb{Q}}
\newcommand{\R}{\mathbb{R}}
\newcommand{\C}{\mathbb{C}}
\newcommand{\F}{\mathbb{F}}
\newcommand{\A}{\mathbb{A}}
\newcommand{\PP}{\mathbb{P}}
\newcommand{\E}{\mathbb{E}}
\newcommand{\eps}{\varepsilon}
\newcommand{\dd}{\,\mathrm{d}}
\begin{document}
\section*{UnsolvedMath ID: 30001215}
\textbf{Source problem code:} catalogue entry
\subsection{Definitions and mathematical statement}
Let \(k\) be a field and let \(A\) be a commutative, unital, finitely generated \(k\)-algebra equipped with a Poisson bracket. Thus \(\{-,-\}:A\times A\to A\) is a \(k\)-bilinear Lie bracket satisfying \(\{a,bc\}=\{a,b\}c+b\{a,c\}\). These are conventions for interpreting the terminology; the supplied record does not specify the base field or additional hypotheses on \(A\).

An ideal \(I\subseteq A\) is a Poisson ideal if \(\{A,I\}\subseteq I\). Write
\[
\operatorname{PSpec}(A)=\{P\subsetneq A:P\text{ is a prime ideal and a Poisson ideal}\},
\]
with the topology induced from \(\operatorname{Spec}(A)\). The Poisson core of an ideal \(J\) is the largest Poisson ideal contained in \(J\), namely
\[
\operatorname{Pcore}(J)=\sum_{\substack{I\subseteq J\\ I\text{ a Poisson ideal}}}I.
\]
For \(P\in\operatorname{PSpec}(A)\), use the following conventions:
\begin{itemize}
\item \(P\) is Poisson-primitive if \(P=\operatorname{Pcore}(\mathfrak m)\) for some maximal ideal \(\mathfrak m\) of \(A\).
\item The bracket on \(A/P\) extends to its fraction field \(K_P=\operatorname{Frac}(A/P)\). Its Poisson center is
\[
Z_{\mathrm{Pois}}(K_P)=\{z\in K_P:\{z,u\}=0\text{ for every }u\in K_P\}.
\]
The ideal \(P\) is Poisson-rational if \(Z_{\mathrm{Pois}}(K_P)\) is algebraic over \(k\).
\item \(P\) is Poisson-locally-closed if \(\{P\}\) is locally closed in \(\operatorname{PSpec}(A)\), meaning it is the intersection of an open subset and a closed subset of that space.
\end{itemize}
The Poisson Dixmier--Moeglin equivalence for \(A\), with these conventions, is the assertion
\[
\forall P\in\operatorname{PSpec}(A),\qquad
P\text{ is Poisson-primitive}
\ \Longleftrightarrow\ 
P\text{ is Poisson-rational}
\ \Longleftrightarrow\ 
P\text{ is Poisson-locally-closed}.
\]
\textbf{Question represented by the record.} Does this equivalence hold for affine Poisson algebras in the intended setting, and under what hypotheses on \(k\) and \(A\) do these three classes coincide? The intended setting is unresolved: the record specifies neither characteristic nor algebraic closedness of \(k\), nor further restrictions on \(A\). Consequently, this draft defines the equivalence precisely for specified \((k,A)\), but does not silently turn the extraction into an unrestricted assertion over all fields and all affine Poisson algebras. The universal validity question and the request to determine sufficient or necessary conditions are distinct targets, although the source presents them together.
\subsection{Short English statement}
For each Poisson prime ideal, compare being the Poisson core of a maximal ideal, having fraction-field Poisson center algebraic over the base field, and being a locally closed point of the Poisson prime spectrum. The problem asks when these conditions are equivalent. The intended field hypotheses and algebra class are missing, so the scope requires human mathematical review.
\subsection{Sources}
\begin{itemize}
\item Supplied local registry record, permanent UnsolvedMath ID \(30001215\), title \emph{Poisson Dixmier--Moeglin Equivalence}, source problem code \texttt{catalogue entry}, in \emph{Oberwolfach Reports Open Problems}. The record gives proposed year 2009 and names no individual proposer.
\item Supplied source citation: \emph{Mini-Workshop: Non-Negativity is a Quantum Phenomenon} (2009), Oberwolfach Reports, pp.~767--800, DOI \texttt{10.4171/owr/2009/14}. \url{https://doi.org/10.4171/owr/2009/14}. This is also the supplied literature-source link labeled ``original OWR report.'' No precise problem page or speaker attribution is supplied.
\item Supplied literature-source label: ``Luo--Wang--Wu primary paper.'' \url{https://arxiv.org/abs/1908.06542}. The record supplies no paper title or specific theorem reference.
\item Literature assessment dated exactly \texttt{2026-08-21}: the supplied curation labels the problem open and states that the equivalence is established for broad classes but remains conditional or open in the unrestricted setting represented by the extraction. This is supplied dated curation, not an independently verified literature assessment. The background explicitly requests independent verification of time-sensitive status and cited claims.
\item The supplied statement status is \texttt{corrected\_verified}; its verification note describes meaning-preserving semantic, typographic, and LaTeX cleanup without unsupported source-dependent additions. This label is source metadata and does not establish mathematical correctness or independent source verification by this draft.
\end{itemize}
\end{document}
